\documentclass[conference]{IEEEtran}
\IEEEoverridecommandlockouts
% The preceding line is only needed to identify funding in the first footnote. If that is unneeded, please comment it out.
\usepackage{cite}
\usepackage{amsmath,amssymb,amsfonts}
\usepackage{graphicx}
\usepackage{textcomp}
\usepackage{xcolor}
\usepackage{booktabs}
\usepackage{multirow}
\usepackage{array}
\usepackage{url}
\usepackage{hyperref}
\usepackage{balance}
\usepackage{svg}
\usepackage{tabularx}
\usepackage{tabularx}
\usepackage{xurl}
\usepackage{seqsplit}
\usepackage{ragged2e}
\def\BibTeX{{\rm B\kern-.05em{\sc i\kern-.025em b}\kern-.08em
    T\kern-.1667em\lower.7ex\hbox{E}\kern-.125emX}}
\begin{document}

\title{Exposing Per-Pod OverlayFS Inode Overhead in Kubernetes via eBPF: Whiteout Amplification, Copy-Up Cost, and Transient Churn\\

\thanks{}
}

\author{Anonymous Author}

\maketitle

\begin{abstract}
In Kubernetes clusters, pods run inside containers whose root filesystems are managed by containerd using OverlayFS as its default snapshotter. The copy-up and whiteout mechanisms of this union filesystem introduce inode overhead that is invisible to conventional node-level monitoring tools, leaving cluster administrators unable to attribute inode pressure to individual pods. This paper presents a dual-layer eBPF framework that simultaneously traces VFS and OverlayFS kernel operations, maps events to Kubernetes pods via cgroup identity, and exports per-pod inode metrics through Prometheus. Controlled evaluation on a MicroK8s cluster reveals two novel findings: lower-layer file deletions produce a consistent 3:1 VFS metadata amplification factor due to whiteout marker creation, while copy-up operations exhibit a 1:1 inode consumption ratio that remains constant regardless of container image layer depth. Comparative analysis confirms that neither standard Kubernetes monitoring tools nor state-of-the-art eBPF security tracers provide equivalent per-pod semantic visibility. The framework operates at under 11 millicores CPU overhead, confirming production viability.
\end{abstract}

\begin{IEEEkeywords}
eBPF, container orchestration, Linux filesystem, storage resource management, cloud-native observability
\end{IEEEkeywords}

\section{Introduction}

In Kubernetes, each pod's root filesystem is assembled by containerd - the default container runtime using OverlayFS as its snapshotter. OverlayFS stacks read-only image layers beneath a writable upper layer, enabling efficient image sharing across pods on the same node. However, this architecture introduces two inode-consuming mechanisms at the kernel level: copy-up operations, which duplicate a lower-layer file into the upper layer when a container modifies it, and whiteout markers, which allocate new upper-layer inodes to mask deletions of read-only lower-layer files. As pod density increases, the cumulative inode overhead from these mechanisms can silently exhaust the host node's inode pool producing filesystem errors that cause scheduling failures regardless of available disk capacity.

Despite the documented risk of inode exhaustion in OverlayFS-based Kubernetes deployments, existing monitoring tools cannot attribute this pressure to individual pods. Standard tools such as cAdvisor and Prometheus Node Exporter provide node-level filesystem metrics but lack per-pod granularity and have no visibility into OverlayFS-internal kernel operations. eBPF-based security tracers observe filesystem activity at the VFS syscall boundary but cannot distinguish the semantic origin of individual VFS operations. Consequently, cluster administrators have no mechanism to identify which pods are generating inode pressure, what OverlayFS operations are responsible, or when node-level exhaustion is approaching.

This paper makes the following contributions:

\begin{itemize}
\item A dual-layer eBPF framework that simultaneously traces both VFS-layer and OverlayFS-layer kernel operations, maps events to Kubernetes pods via cgroup identity, and exports per-pod inode metrics through Prometheus.
\item Empirically quantify for the first time that lower-layer file deletions produce a consistent $3\!:\!1$ VFS metadata amplification factor due to whiteout marker creation.
\item Empirically establish that copy-up operations consume inodes at a precise $1\!:\!1$ ratio per operation.
\item Demonstrate through comparative evaluation that no existing monitoring tool provides equivalent per-pod OverlayFS semantic visibility.
\end{itemize}

The remainder of this paper is organized as follows. Section II describes the framework design and experimental setup. Section III presents the evaluation results. Section IV discusses the findings and their implications. Section V concludes the paper.

\section{System Design}
\label{sec:design}

The proposed framework consists of two cooperating components: a kernel-resident eBPF tracing module written in C and a user-space aggregation daemon written in Go. The framework is deployed as a privileged DaemonSet on each Kubernetes worker node. Fig.~\ref{fig:architecture} illustrates the end-to-end architecture.

% ----- Architecture Figure -----
\begin{figure}[!htbp]
    \centering
    \includesvg[width=0.95\columnwidth]{Kernal Tracing.drawio}
    \caption{End-to-end framework architecture.}
    \label{fig:architecture}
\end{figure}

\subsection{Dual-Layer eBPF Probe Architecture}
\label{subsec:probes}

The eBPF module attaches probes at two distinct kernel layers to capture both physical inode allocations and overlay-specific semantic operations.

\subsubsection{VFS Layer (fentry Probes)}
Function-entry (\texttt{fentry}) probes are attached to core Virtual File System routines using BPF Type Format (BTF) metadata. The \texttt{fentry} mechanism patches the first instruction of the target function to redirect execution into the eBPF handler with near-zero overhead (no \texttt{int3} interrupt or stack-frame copying), receiving fully typed function arguments. Table~\ref{tab:vfs_probes} lists the traced VFS functions and their inode impact.

\begin{table}[t]
\centering
\caption{VFS-layer probes attached via \texttt{fentry}.}
\label{tab:vfs_probes}
\small
\resizebox{\columnwidth}{!}{
\begin{tabular}{llcl}
\toprule
\textbf{Kernel Function} & \textbf{Probe} & \textbf{$\Delta$ Inode} & \textbf{Purpose} \\
\midrule
\texttt{vfs\_create}  & fentry & $+1$ & File creation \\
\texttt{vfs\_mkdir}   & fentry & $+1$ & Directory creation \\
\texttt{vfs\_symlink} & fentry & $+1$ & Symlink creation \\
\texttt{vfs\_unlink}  & fentry & $-1$ & File deletion \\
\texttt{vfs\_rmdir}   & fentry & $-1$ & Directory removal \\
\texttt{vfs\_rename}  & fentry & $0$  & Rename (no new inode) \\
\texttt{vfs\_link}    & fentry & $0$  & Hard link (no new inode) \\
\texttt{vfs\_mknod}   & kprobe & $+1$ & Device node creation \\
\bottomrule
\end{tabular}
}
\end{table}

\subsubsection{OverlayFS Layer (kprobes)}
Kernel probes (\texttt{kprobes}) are bound to internal OverlayFS routines, which are not always BTF-annotated and whose symbol names may vary across kernel versions. The daemon iterates through a candidate symbol list per probe and attaches to the first resolved match, ensuring cross-version compatibility. Table~\ref{tab:ovl_probes} lists the traced OverlayFS functions.

\begin{table}[t]
\centering
\caption{OverlayFS-layer probes attached via \texttt{kprobe}.}
\label{tab:ovl_probes}
\small
\resizebox{\columnwidth}{!}{
\begin{tabular}{lcl}
\toprule
\textbf{Kernel Function(s)} & \textbf{$\Delta$ Inode} & \textbf{Purpose} \\
\midrule
\texttt{ovl\_copy\_up\_one} / \texttt{ovl\_copy\_up\_start} & $+1$ & Copy-up: lower to upper layer \\
\midrule
\texttt{ovl\_unlink} / \texttt{ovl\_do\_unlink} / \texttt{ovl\_remove} & $-1$ & OverlayFS-level file deletion \\
\midrule
\texttt{ovl\_cleanup\_and\_whiteout} & $0$ & Whiteout marker creation \\
\texttt{ovl\_lookup} & $0$ & Layer traversal resolution \\
\bottomrule
\end{tabular}
}
\end{table}

\subsection{Cgroup-Based Event Filtering}
\label{subsec:filtering}

Every probe begins execution with a two-stage filtering sequence. First, the BPF helper \texttt{bpf\_get\_current\_cgroup\_id()} retrieves the cgroup~v2 identifier of the process triggering the syscall. This identifier is looked up in a BPF hash map (\texttt{cgroup\_filter}) that serves as a whitelist of monitored container cgroups, populated and periodically refreshed by the user-space daemon. If the calling process does not belong to a monitored pod, the probe returns immediately with negligible overhead. Second, if the cgroup is whitelisted, a slot is reserved in a lock-free BPF ring buffer (4\,MB capacity). The probe then populates a fixed-size event structure comprising the operation type identifier, the cgroup ID, and an inode-delta value ($+1$, $-1$, or $0$) before asynchronously submitting the event to user space. The ring buffer was chosen over the legacy \texttt{PERF\_EVENT\_ARRAY} for its superior memory efficiency and absence of per-CPU buffer fragmentation.


\subsection{User-Space Daemon and Metric Export}
\label{subsec:daemon}

The Go-based daemon executes four principal functions:

\textbf{(1)~eBPF Lifecycle Management.} The daemon loads the pre-compiled eBPF object into the kernel, validates it via the kernel verifier, and attaches all probes. For \texttt{fentry} probes, BTF-based binding is used; for \texttt{kprobes}, the daemon tests multiple candidate symbol names per function to ensure compatibility across kernel versions.

\textbf{(2)~Cgroup-to-Pod Resolution.} Two concurrent mapping tables translate raw kernel cgroup IDs into Kubernetes metadata. A \emph{cgroup inode map} is populated by periodically walking \texttt{/sys/fs/cgroup/kubepods.slice} and extracting the filesystem inode number of each cgroup directory (identical to the cgroup IDs reported in kernel space). A \emph{pod identity map} is maintained via a persistent Kubernetes API watch stream, mapping each pod's UID to its display name and namespace. Synchronization occurs on a 5-second interval, with a fallback full API re-query for newly scheduled pods.

\textbf{(3)~Event Consumption Pipeline.} Events from the BPF ring buffer are consumed into a high-capacity Go channel (1\,M elements). Each event undergoes binary decoding, cgroup path resolution, pod UID extraction, and pod metadata association, producing a fully attributed event record.

\textbf{(4)~Prometheus Metric Exportation.} Three metrics are exported via an HTTP endpoint, as summarized in Table~\ref{tab:metrics}. Upon pod deletion, associated label sets are cleaned up to prevent stale time series.

\begin{table}[t]
\centering
\caption{Prometheus metrics exported by the framework.}
\label{tab:metrics}
\small

\renewcommand{\arraystretch}{1.2}

\begin{tabularx}{\columnwidth}{ 
>{\RaggedRight\arraybackslash}p{4cm}
c
>{\RaggedRight\arraybackslash}X
}

\toprule
\textbf{Metric} & \textbf{Type} & \textbf{Purpose} \\
\midrule

\texttt{\seqsplit{\detokenize{container_inode_operation_total}}} & Counter & Per-pod, per-operation cumulative count \\

\texttt{\seqsplit{\detokenize{container_inode_net_delta}}} & Gauge & Running net inode $\Delta$ per pod \\

\texttt{\seqsplit{\detokenize{node_inode_free}}} & Gauge & Host free inode count \\

\bottomrule
\end{tabularx}

\end{table}

\section{Prepare Your Paper Before Styling}
Before you begin to format your paper, first write and save the content as a 
separate text file. Complete all content and organizational editing before 
formatting. Please note sections \ref{AA}--\ref{SCM} below for more information on 
proofreading, spelling and grammar.

Keep your text and graphic files separate until after the text has been 
formatted and styled. Do not number text heads---{\LaTeX} will do that 
for you.

\subsection{Abbreviations and Acronyms}\label{AA}
Define abbreviations and acronyms the first time they are used in the text, 
even after they have been defined in the abstract. Abbreviations such as 
IEEE, SI, MKS, CGS, ac, dc, and rms do not have to be defined. Do not use 
abbreviations in the title or heads unless they are unavoidable.

\subsection{Units}
\begin{itemize}
\item Use either SI (MKS) or CGS as primary units. (SI units are encouraged.) English units may be used as secondary units (in parentheses). An exception would be the use of English units as identifiers in trade, such as ``3.5-inch disk drive''.
\item Avoid combining SI and CGS units, such as current in amperes and magnetic field in oersteds. This often leads to confusion because equations do not balance dimensionally. If you must use mixed units, clearly state the units for each quantity that you use in an equation.
\item Do not mix complete spellings and abbreviations of units: ``Wb/m\textsuperscript{2}'' or ``webers per square meter'', not ``webers/m\textsuperscript{2}''. Spell out units when they appear in text: ``. . . a few henries'', not ``. . . a few H''.
\item Use a zero before decimal points: ``0.25'', not ``.25''. Use ``cm\textsuperscript{3}'', not ``cc''.)
\end{itemize}

\subsection{Equations}
Number equations consecutively. To make your 
equations more compact, you may use the solidus (~/~), the exp function, or 
appropriate exponents. Italicize Roman symbols for quantities and variables, 
but not Greek symbols. Use a long dash rather than a hyphen for a minus 
sign. Punctuate equations with commas or periods when they are part of a 
sentence, as in:
\begin{equation}
a+b=\gamma\label{eq}
\end{equation}

Be sure that the 
symbols in your equation have been defined before or immediately following 
the equation. Use ``\eqref{eq}'', not ``Eq.~\eqref{eq}'' or ``equation \eqref{eq}'', except at 
the beginning of a sentence: ``Equation \eqref{eq} is . . .''

\subsection{\LaTeX-Specific Advice}

Please use ``soft'' (e.g., \verb|\eqref{Eq}|) cross references instead
of ``hard'' references (e.g., \verb|(1)|). That will make it possible
to combine sections, add equations, or change the order of figures or
citations without having to go through the file line by line.

Please don't use the \verb|{eqnarray}| equation environment. Use
\verb|{align}| or \verb|{IEEEeqnarray}| instead. The \verb|{eqnarray}|
environment leaves unsightly spaces around relation symbols.

Please note that the \verb|{subequations}| environment in {\LaTeX}
will increment the main equation counter even when there are no
equation numbers displayed. If you forget that, you might write an
article in which the equation numbers skip from (17) to (20), causing
the copy editors to wonder if you've discovered a new method of
counting.

{\BibTeX} does not work by magic. It doesn't get the bibliographic
data from thin air but from .bib files. If you use {\BibTeX} to produce a
bibliography you must send the .bib files. 

{\LaTeX} can't read your mind. If you assign the same label to a
subsubsection and a table, you might find that Table I has been cross
referenced as Table IV-B3. 

{\LaTeX} does not have precognitive abilities. If you put a
\verb|\label| command before the command that updates the counter it's
supposed to be using, the label will pick up the last counter to be
cross referenced instead. In particular, a \verb|\label| command
should not go before the caption of a figure or a table.

Do not use \verb|\nonumber| inside the \verb|{array}| environment. It
will not stop equation numbers inside \verb|{array}| (there won't be
any anyway) and it might stop a wanted equation number in the
surrounding equation.

\subsection{Some Common Mistakes}\label{SCM}
\begin{itemize}
\item The word ``data'' is plural, not singular.
\item The subscript for the permeability of vacuum $\mu_{0}$, and other common scientific constants, is zero with subscript formatting, not a lowercase letter ``o''.
\item In American English, commas, semicolons, periods, question and exclamation marks are located within quotation marks only when a complete thought or name is cited, such as a title or full quotation. When quotation marks are used, instead of a bold or italic typeface, to highlight a word or phrase, punctuation should appear outside of the quotation marks. A parenthetical phrase or statement at the end of a sentence is punctuated outside of the closing parenthesis (like this). (A parenthetical sentence is punctuated within the parentheses.)
\item A graph within a graph is an ``inset'', not an ``insert''. The word alternatively is preferred to the word ``alternately'' (unless you really mean something that alternates).
\item Do not use the word ``essentially'' to mean ``approximately'' or ``effectively''.
\item In your paper title, if the words ``that uses'' can accurately replace the word ``using'', capitalize the ``u''; if not, keep using lower-cased.
\item Be aware of the different meanings of the homophones ``affect'' and ``effect'', ``complement'' and ``compliment'', ``discreet'' and ``discrete'', ``principal'' and ``principle''.
\item Do not confuse ``imply'' and ``infer''.
\item The prefix ``non'' is not a word; it should be joined to the word it modifies, usually without a hyphen.
\item There is no period after the ``et'' in the Latin abbreviation ``et al.''.
\item The abbreviation ``i.e.'' means ``that is'', and the abbreviation ``e.g.'' means ``for example''.
\end{itemize}
An excellent style manual for science writers is \cite{b7}.

\subsection{Authors and Affiliations}
\textbf{The class file is designed for, but not limited to, six authors.} A 
minimum of one author is required for all conference articles. Author names 
should be listed starting from left to right and then moving down to the 
next line. This is the author sequence that will be used in future citations 
and by indexing services. Names should not be listed in columns nor group by 
affiliation. Please keep your affiliations as succinct as possible (for 
example, do not differentiate among departments of the same organization).

\subsection{Identify the Headings}
Headings, or heads, are organizational devices that guide the reader through 
your paper. There are two types: component heads and text heads.

Component heads identify the different components of your paper and are not 
topically subordinate to each other. Examples include Acknowledgments and 
References and, for these, the correct style to use is ``Heading 5''. Use 
``figure caption'' for your Figure captions, and ``table head'' for your 
table title. Run-in heads, such as ``Abstract'', will require you to apply a 
style (in this case, italic) in addition to the style provided by the drop 
down menu to differentiate the head from the text.

Text heads organize the topics on a relational, hierarchical basis. For 
example, the paper title is the primary text head because all subsequent 
material relates and elaborates on this one topic. If there are two or more 
sub-topics, the next level head (uppercase Roman numerals) should be used 
and, conversely, if there are not at least two sub-topics, then no subheads 
should be introduced.

\subsection{Figures and Tables}
\paragraph{Positioning Figures and Tables} Place figures and tables at the top and 
bottom of columns. Avoid placing them in the middle of columns. Large 
figures and tables may span across both columns. Figure captions should be 
below the figures; table heads should appear above the tables. Insert 
figures and tables after they are cited in the text. Use the abbreviation 
``Fig.~\ref{fig}'', even at the beginning of a sentence.

\begin{table}[htbp]
\caption{Table Type Styles}
\begin{center}
\begin{tabular}{|c|c|c|c|}
\hline
\textbf{Table}&\multicolumn{3}{|c|}{\textbf{Table Column Head}} \\
\cline{2-4} 
\textbf{Head} & \textbf{\textit{Table column subhead}}& \textbf{\textit{Subhead}}& \textbf{\textit{Subhead}} \\
\hline
copy& More table copy$^{\mathrm{a}}$& &  \\
\hline
\multicolumn{4}{l}{$^{\mathrm{a}}$Sample of a Table footnote.}
\end{tabular}
\label{tab1}
\end{center}
\end{table}

\begin{figure}[htbp]
\centerline{\includegraphics{fig1.png}}
\caption{Example of a figure caption.}
\label{fig}
\end{figure}

Figure Labels: Use 8 point Times New Roman for Figure labels. Use words 
rather than symbols or abbreviations when writing Figure axis labels to 
avoid confusing the reader. As an example, write the quantity 
``Magnetization'', or ``Magnetization, M'', not just ``M''. If including 
units in the label, present them within parentheses. Do not label axes only 
with units. In the example, write ``Magnetization (A/m)'' or ``Magnetization 
\{A[m(1)]\}'', not just ``A/m''. Do not label axes with a ratio of 
quantities and units. For example, write ``Temperature (K)'', not 
``Temperature/K''.

\section*{Acknowledgment}

The preferred spelling of the word ``acknowledgment'' in America is without 
an ``e'' after the ``g''. Avoid the stilted expression ``one of us (R. B. 
G.) thanks $\ldots$''. Instead, try ``R. B. G. thanks$\ldots$''. Put sponsor 
acknowledgments in the unnumbered footnote on the first page.

\section*{References}

Please number citations consecutively within brackets \cite{b1}. The 
sentence punctuation follows the bracket \cite{b2}. Refer simply to the reference 
number, as in \cite{b3}---do not use ``Ref. \cite{b3}'' or ``reference \cite{b3}'' except at 
the beginning of a sentence: ``Reference \cite{b3} was the first $\ldots$''

Number footnotes separately in superscripts. Place the actual footnote at 
the bottom of the column in which it was cited. Do not put footnotes in the 
abstract or reference list. Use letters for table footnotes.

Unless there are six authors or more give all authors' names; do not use 
``et al.''. Papers that have not been published, even if they have been 
submitted for publication, should be cited as ``unpublished'' \cite{b4}. Papers 
that have been accepted for publication should be cited as ``in press'' \cite{b5}. 
Capitalize only the first word in a paper title, except for proper nouns and 
element symbols.

For papers published in translation journals, please give the English 
citation first, followed by the original foreign-language citation \cite{b6}.

\begin{thebibliography}{00}
\bibitem{b1} G. Eason, B. Noble, and I. N. Sneddon, ``On certain integrals of Lipschitz-Hankel type involving products of Bessel functions,'' Phil. Trans. Roy. Soc. London, vol. A247, pp. 529--551, April 1955.
\bibitem{b2} J. Clerk Maxwell, A Treatise on Electricity and Magnetism, 3rd ed., vol. 2. Oxford: Clarendon, 1892, pp.68--73.
\bibitem{b3} I. S. Jacobs and C. P. Bean, ``Fine particles, thin films and exchange anisotropy,'' in Magnetism, vol. III, G. T. Rado and H. Suhl, Eds. New York: Academic, 1963, pp. 271--350.
\bibitem{b4} K. Elissa, ``Title of paper if known,'' unpublished.
\bibitem{b5} R. Nicole, ``Title of paper with only first word capitalized,'' J. Name Stand. Abbrev., in press.
\bibitem{b6} Y. Yorozu, M. Hirano, K. Oka, and Y. Tagawa, ``Electron spectroscopy studies on magneto-optical media and plastic substrate interface,'' IEEE Transl. J. Magn. Japan, vol. 2, pp. 740--741, August 1987 [Digests 9th Annual Conf. Magnetics Japan, p. 301, 1982].
\bibitem{b7} M. Young, The Technical Writer's Handbook. Mill Valley, CA: University Science, 1989.
\end{thebibliography}
\vspace{12pt}
\color{red}
IEEE conference templates contain guidance text for composing and formatting conference papers. Please ensure that all template text is removed from your conference paper prior to submission to the conference. Failure to remove the template text from your paper may result in your paper not being published.

\end{document}
